\chapter{Defence and security as engineered systems}
\label{vol12:ch09}

\epigraph{Defence engineering is engineering: the constraints are stricter, the failures are more public, the ethics are sharper, the responsibility is the same.}{Author}

\chapmeta{Half-life: medium-life. Archetypes: control, ethics, failure. Exercise target: 22. Project track: design (a small defence-relevant component with dual-use analysis). Hazard class: high (existential and civilizational stakes). Reader path default: standard, with the defence-as-domain and dual-use sections core and the procurement and verification sections mastery. Named cases: Patriot timing software failure Dhahran 1991 (\cite{acc:gao-patriot-1992}); Stuxnet 2010 (\cite{paper:langner-stuxnet-2013}); F-35 ALIS programme (\cite{acc:gao-f35-alis-2020}).}

\noindent
Defence engineering is included in this volume because the project's wide stance does not permit it to be omitted: defence procurement, military command-and-control, cyber defence, critical-infrastructure protection, and arms-control verification are all civilization-scale engineering domains, and the engineering profession bears continuing responsibility for the dual-use technologies it produces. The chapter is openly normative about this responsibility while remaining technically descriptive about the engineering substance. Global defence spending in 2024 was approximately \$2.4 trillion (SIPRI), of which a substantial fraction supports defence engineering and the defence-industrial workforce; the engineering profession's involvement is unavoidable.

\section{Defence as engineered domain}\pathtag{core}

Defence engineering covers the full engineering portfolio: structural (warships, aircraft, ground vehicles); propulsion (gas turbines, nuclear, electric, conventional); weapons systems (guns, missiles, electronic countermeasures); sensors (radar, sonar, electro-optical, signals intelligence); communications (encrypted, low-probability-of-intercept, beyond-line-of-sight); command-and-control (decision support; doctrinal; AI-supported); logistics (maintenance; supply chain; basing); cyber (offensive and defensive); space (Volume XII Chapter 10). The engineering disciplines involved overlap substantially with the civilian-engineering portfolio with the addition of survivability against adversarial action, the operational regime under conflict conditions, and the procurement and acquisition regulatory regime.

The contemporary defence-engineering trajectory is increasingly software-and-electronics dominated. A modern combat aircraft is approximately $30\%$ software-and-electronics by cost (an F-35 has approximately $9 \times 10^6$ lines of code in onboard software); a warship is comparable; a ground vehicle is increasingly so. The implication: the defence-engineering profession is increasingly a software-engineering profession, with the software-engineering practices, quality challenges, and failure modes that Volume X Chapter 9 documented. The Patriot Dhahran software failure (\cite{acc:gao-patriot-1992}) and the F-35 ALIS programme difficulties (\cite{acc:gao-f35-alis-2020}) are both fundamentally software-engineering failures within larger systems.

\begin{archetype}[Control: defence systems]
A defence system combines sensing (detection of threat), command (decision among engagement options), control (execution of engagement), and effects (kinetic or non-kinetic). The engineering scope: the closed-loop with all four stages operating under adversarial conditions, with the adversary actively trying to defeat detection, command, control, and effects. The contemporary trend toward autonomous and human-on-the-loop systems compresses the time scale, raises the technology stakes, and introduces ethical questions that the discipline is still working out. The Vincennes shootdown of Iran Air 655 in 1988 illustrated the operational consequences of complex automated combat-decision systems under stress; the contemporary debate around autonomous weapons (lethal autonomous weapons systems; LAWS) is the continuing engineering-and-ethical question.
\end{archetype}

\section{Defence-industrial supply chains and procurement}\pathtag{mastery}

The defence-industrial supply chain is among the most complex engineered supply chains in modern industry. A single combat aircraft programme involves thousands of suppliers across dozens of countries; the F-35 programme has approximately $1{,}400$ suppliers across $9$ partner nations and many more secondary suppliers. The procurement regulatory regime (FAR/DFAR in the US; equivalent in other countries) imposes substantial overhead but also provides accountability for the cost and performance of multi-decade programmes.

The engineering implications of defence procurement: long programme lifecycles (typically $20$-$40$ years from concept to deployment for major systems; another $20$-$50$ years of operational service); the multi-generation requirements management; the integration of legacy and new systems; the configuration management at fleet scale; the post-deployment support engineering. The procurement and acquisition profession is itself an engineering discipline that combines technical, financial, and institutional considerations.

The F-35 ALIS programme (\cite{acc:gao-f35-alis-2020}) is a representative procurement-and-software case: the Autonomic Logistics Information System, designed as a comprehensive logistics-management software for the aircraft fleet, was substantially over budget, late, and operationally problematic across more than a decade of development; the eventual replacement (ODIN, Operational Data Integrated Network) is a fresh-start effort whose own success remains to be demonstrated. The lesson: defence software programmes face the same engineering challenges as commercial software programmes, with the additional difficulty of the procurement regulatory regime, the multi-generation requirements drift, and the operational-environment constraints that make iteration difficult.

\section{Sensors, networks, and command-and-control}\pathtag{standard}

The contemporary defence sensor portfolio includes: radar (search, fire-control, tracking; ground, ship, air, space); sonar (active, passive; surface, submarine); electro-optical and infrared (imaging, targeting, missile-warning); signals intelligence (communications interception; electronic-warfare emitter identification); space-based imagery (electro-optical, synthetic-aperture radar, hyperspectral); ground-moving-target indication; persistent surveillance (long-endurance UAV; aerostat; satellite constellations). The fused sensor picture is the input to the command-and-control system.

The contemporary command-and-control system trend is the network-centric warfare concept: distributed sensors feeding a common operational picture; geographically dispersed decision-makers operating against the same data; effects platforms tasked dynamically across organisational boundaries. The engineering challenge is the network: the bandwidth and latency requirements, the survivability of the communications under adversarial conditions, the encryption and authentication that prevents adversary access. The engineering reality has lagged the doctrinal aspiration in many programmes; the contemporary work on JADC2 (Joint All-Domain Command and Control) in the US is the latest attempt to deliver on the network-centric vision.

\section{Cyber defence at the state scale}\pathtag{standard}

Cyber defence is the engineering of state-scale protection against adversarial cyber operations. The threat portfolio includes: ransomware against critical infrastructure (Colonial Pipeline 2021; HSE Ireland 2021); espionage against government and industry (SolarWinds 2020; Microsoft Exchange 2021); destructive attacks against industrial-control systems (Stuxnet 2010; \cite{paper:langner-stuxnet-2013}; NotPetya 2017 with the substantial collateral damage to civilian shipping and logistics); disinformation operations against electoral and political processes (Russian operations against US elections 2016 onward; cross-jurisdictional). The defensive engineering combines: detection (network monitoring; endpoint detection-and-response; intelligence-sharing); prevention (patching; configuration management; identity-and-access; supply-chain risk management); response (incident response; recovery engineering; forensic analysis); deterrence (offensive cyber operations; attribution; sanctions; the contested deterrence theory).

The Stuxnet operation \cite{paper:langner-stuxnet-2013} of 2009-2010 is the contemporary canonical case: a sophisticated worm targeting Siemens industrial-control systems specifically configured for the Iranian nuclear-enrichment programme; the operation produced substantial damage to the centrifuge cascade while remaining undetected for an extended period; the eventual discovery and public attribution shaped the subsequent state-cyber-operations doctrine and the recognition by the civilian engineering profession of the industrial-control-system cyber-vulnerability portfolio. The continuing engineering response: the IEC 62443 and equivalent standards for industrial-control-system cybersecurity; the OT (operational technology) cybersecurity profession that emerged in the post-Stuxnet decade; the dedicated CISA, ENISA, and equivalent agencies in major jurisdictions.

\section{Critical-infrastructure defence}\pathtag{core}

Critical-infrastructure defence is the protection of the civilian infrastructures the previous chapters of this volume have described: energy, water, transport, food, healthcare, communications. The threat portfolio combines kinetic (terrorism; sabotage; state-level military action against infrastructure as in the Russia-Ukraine war 2022 onward) and non-kinetic (cyber; supply-chain compromise; insider threat; electromagnetic-pulse). The engineering response is layered: hardening (the physical security; cyber security; access control; supply-chain integrity); redundancy and resilience (the engineered ability to maintain function despite partial loss of components); recovery (the engineered ability to restore function after attack); deterrence (through attribution and consequence).

The contemporary engineering profession's contribution: the explicit threat modelling of critical infrastructure (typically performed under classified or sensitive-but-unclassified frameworks; some publicly published as Cybersecurity \& Infrastructure Security Agency (CISA) sector-specific guidance and equivalent); the design-against-attack practices in new infrastructure development; the post-incident forensic-and-recovery engineering. The engineering and policy interface: the determination of what to protect, against what threat, to what level of effort, with what investment, is a joint engineering-and-political decision.

\section{Dual-use technology and export control}\pathtag{core}

Dual-use technology is the engineering output that has both civilian and military applications. The major dual-use domains: nuclear (energy and weapons); biotechnology (medicine and biological weapons); chemical (industrial chemistry and chemical weapons); aerospace (commercial and military aviation; satellites; missiles); cyber (defensive and offensive); AI (commercial and military applications); materials (semiconductor manufacturing equipment; high-strength alloys; cryptographic systems). The export-control regimes (the Nuclear Suppliers Group; the Australia Group for chemical and biological; the Missile Technology Control Regime; the Wassenaar Arrangement for conventional and dual-use; the EU dual-use regulation; the US ITAR and EAR) impose restrictions on the export of dual-use items to specified jurisdictions and end-users.

The engineering profession's responsibility includes the assessment of dual-use risk for engineering outputs and the engagement with the export-control regime as part of the engineering scope. The contemporary debates: the appropriate scope of export control for semiconductor manufacturing equipment (the US 2022 expanded controls on advanced semiconductor exports to China); the appropriate scope of export control for cyber-intrusion tools (the 2021 Wassenaar additions; the continuing implementation challenges); the appropriate scope of export control for AI capabilities. Engineering judgements about dual-use are not optional; the choice to proceed with engineering work on a dual-use technology is a choice about its proliferation potential and its end-uses, and the engineering profession is one of the actors whose judgement matters.

\begin{principle}[Dual-use as engineering responsibility]
Any engineering technology with potential dual-use application carries an associated responsibility for the engineering profession to consider its proliferation potential, its end-use implications, and its export-control treatment. The responsibility is not transferred to the procurement or policy functions; it is shared with them. The engineering profession's failure to engage with the dual-use question for its work outputs is one of the failure modes the chapter's later sections document.
\end{principle}

\section{Arms control and verification engineering}\pathtag{mastery}

Arms-control verification is the engineering of mechanisms by which compliance with treaty obligations can be observed and confirmed. The major examples: the IAEA safeguards regime for nuclear non-proliferation (the comprehensive safeguards under the NPT; the additional protocol for enhanced inspection); the CTBT verification regime for nuclear-test detection (the International Monitoring System with seismic, hydroacoustic, infrasound, and radionuclide sensors; the on-site inspection regime); the OPCW verification regime for chemical-weapons (the chemical-industry declarations and inspections; the destruction verification); the Open Skies treaty (unarmed aerial reconnaissance flights over participating-state territory; substantially undermined by US and Russian withdrawal in 2020-2021); the strategic-arms verification under New START (national-technical-means with notification protocols).

The engineering substance: the sensor technology that supports detection; the data-fusion-and-analysis infrastructure; the inspection protocols and the support equipment for inspections; the technical analysis of compliance-or-non-compliance evidence. The verification engineering is a substantial subdiscipline with its own professional community (the Federation of American Scientists; the Bulletin of the Atomic Scientists; the International Network of Engineers and Scientists for Global Responsibility) that the engineering profession has historically supported. The contemporary trajectory is mixed: some arms-control regimes are eroding (Open Skies; the INF Treaty); others are persisting (the IAEA safeguards; the OPCW); new arms-control questions are emerging (autonomous weapons; cyber; space). The engineering profession's continuing contribution to verification is one of the long-standing examples of the profession engaging with civilization-scale ethical responsibilities.

\section{Failure: defence-systems failures}\pathtag{core}

The chapter's failure cases each illustrate a distinct failure mode of defence-engineering.

The Patriot timing software failure at Dhahran on 25 February 1991 \cite{acc:gao-patriot-1992} killed $28$ US service members at a barracks struck by a Scud missile that the Patriot battery failed to engage. The technical mechanism: a software clock-drift accumulation in the system's 24-bit fixed-point time representation; after approximately $100$ hours of continuous operation the drift was sufficient to cause the radar's range-gate to miss the incoming missile. A software patch had been written but not yet deployed to all theatre Patriots; the Dhahran battery had not received it. The engineering lesson: the production-deployment of even minor software patches in critical systems has consequences when not executed timely; the software-quality regime for safety-critical defence systems must include the patch-deployment infrastructure as part of the system.

The Bradley Fighting Vehicle programme is the classical post-Vietnam US defence-procurement-difficulty case: a $20$-year development from concept (1963) to deployment (1981) with substantial requirements drift, cost growth, and operational concerns; documented in the contemporary critique \cite{paper:burton-pentagon-wars-1993}. The eventual operational performance was substantially better than the critique anticipated; the lesson is more about the procurement-process failure modes than about the engineering output. The contemporary US defence-procurement reform efforts (the Adaptive Acquisition Framework introduced in 2020) attempt to address the same patterns.

The F-35 ALIS programme \cite{acc:gao-f35-alis-2020} is the contemporary defence-software case: more than a decade of attempted delivery of a comprehensive logistics-management software; chronic operational deficiencies, cost growth, schedule slippage; eventual decision to replace with a fresh-start effort (ODIN). The lesson: the chronic software-programme failure modes that the commercial software industry has worked to address (Volume X Chapter 9) apply with full force to defence software, with the additional constraints of the defence-procurement regulatory regime, the multi-decade requirements timeline, and the operational-environment difficulty.

\begin{principle}[Defence engineering carries the same discipline as civilian engineering, with stricter constraints and sharper ethics]
The defence-engineering domain operates under the same engineering discipline as the civilian-engineering work that the previous volumes have described, with three differences: the operational regime is more demanding (adversarial conditions; operational stress; performance margins); the procurement and acquisition regime imposes substantial process overhead; the ethical responsibility for dual-use, end-use, and proliferation is sharper. The engineering profession's discipline applies to defence work without remainder, and the engineering profession's ethical responsibility extends to engagement with the dual-use, arms-control, and end-use questions that the work raises.
\end{principle}

\section{Project}\pathtag{core}

\begin{project}[Defence-relevant component with dual-use analysis]
Choose a small defence-relevant engineering component (a drone-detection-and-jamming device; a secure radio for civil-society communications; a supply-chain-provenance tool for hardware integrity; a cryptographic system for state-grade communications). Design the component to a baseline specification. Document the engineering trade-offs (cost vs performance; size vs capability; security vs accessibility); document the dual-use considerations (the civilian applications; the military applications; the misuse potential); document the export-control treatment under the relevant regime; document the ethical reasoning behind the decision to proceed with the engineering work and the conditions under which the work should be paused or modified. The project's deliverable is the technical design plus the dual-use analysis; the analysis is the part of the engineering scope that this chapter is most insistent on.
\end{project}

\section{Exercises}

\subsection*{Calculation}\pathtag{core}

\begin{exercise}
A radar system's clock drift is $\Delta t/t = 10^{-7}$ per second. Compute the cumulative drift over $100$ hours of continuous operation and the effective range-gate offset at the speed of light.
\end{exercise}

\begin{exercise}
A defence-industrial supplier base of $1{,}400$ first-tier suppliers across $9$ countries is producing $200$ aircraft/yr at $\$80$ million unit cost. Compute the annual programme value and the supplier-base economic distribution.
\end{exercise}

\begin{exercise}
A critical-infrastructure cyber-incident-response programme aims to detect, contain, and recover from a major attack within $24$ hours. Compute the staffing required for $24/7$ coverage at $3$-person-shift cycles with overlap.
\end{exercise}

\begin{exercise}
A satellite-based reconnaissance constellation has $100$ satellites at $500\,\text{km}$ altitude with $30$-minute revisit time over a target. Compute the orbital-mechanics constraints (number of orbital planes; satellites per plane) required to achieve the revisit target.
\end{exercise}

\subsection*{Derivation}\pathtag{standard}

\begin{exercise}
Derive the radar range equation for a search radar against a fluctuating target. Identify the key trade-offs among transmitted power, antenna gain, target RCS, and detection threshold.
\end{exercise}

\begin{exercise}
Derive the optimal allocation of cyber-defence investment across detection, prevention, response, and recovery functions given a budget constraint and an adversarial threat model.
\end{exercise}

\begin{exercise}
Derive the bandwidth-and-latency requirements for a tactical network supporting distributed sensor-shooter coordination across $N$ platforms at given engagement timescales.
\end{exercise}

\begin{exercise}
Derive the verification probability for an arms-control regime as a function of the sensor coverage, the false-positive rate, and the detection threshold.
\end{exercise}

\subsection*{Estimation}\pathtag{standard}

\begin{exercise}
Estimate the global defence-engineering workforce and its distribution by domain (aerospace; ground; naval; cyber; space; munitions).
\end{exercise}

\begin{exercise}
Estimate the lead time and capital investment to bring a country to indigenous capability in a critical defence-technology area (e.g., advanced semiconductor manufacturing; jet-engine production).
\end{exercise}

\begin{exercise}
Estimate the global ransomware-against-critical-infrastructure exposure (annual incidents; economic impact; engineering and operational response).
\end{exercise}

\begin{exercise}
Estimate the technical and institutional resource required to verify a comprehensive ban on autonomous-weapons-systems with given technical capabilities and given state participation.
\end{exercise}

\subsection*{Simulation}\pathtag{standard}

\begin{exercise}
Simulate a critical-infrastructure cyber-attack scenario with multi-stage propagation. Identify the defensive configurations that contain the attack versus those that allow propagation.
\end{exercise}

\begin{exercise}
Simulate a defence-procurement programme with realistic requirements-drift, cost-growth, and schedule-slip dynamics. Identify the management and procurement-regime interventions that reduce the failure rate.
\end{exercise}

\begin{exercise}
Simulate an arms-control verification regime under realistic detection and false-alarm parameters. Identify the regime configurations that produce credible compliance assurance.
\end{exercise}

\subsection*{Design}\pathtag{standard}

\begin{exercise}
Design a cyber-defence architecture for an electric-grid operator. Specify the IT/OT segmentation, the detection-and-monitoring infrastructure, the incident-response capability, the recovery engineering.
\end{exercise}

\begin{exercise}
Design a dual-use risk-assessment process for an engineering organisation producing technologies with potential military applications. Specify the criteria, the decision points, the documentation, the institutional accountability.
\end{exercise}

\begin{exercise}
Design a supply-chain-integrity programme for a defence-electronics manufacturer to prevent counterfeit-component infiltration. Specify the supplier-qualification, the inspection, the traceability, the audit programme.
\end{exercise}

\subsection*{Judgement}\pathtag{mastery}

\begin{exercise}
Discuss the Patriot Dhahran software failure as a defence-software engineering failure. Identify the engineering, organisational, and operational decisions that contributed; identify the structural improvements that would prevent recurrence.
\end{exercise}

\begin{exercise}
Argue: the engineering profession has a continuing responsibility to engage with arms-control verification, even when state-level arms-control regimes are eroding. Argue the opposite. Identify the implications for engineering-professional practice.
\end{exercise}

\begin{exercise}
Discuss the appropriate engineering profession's stance on autonomous-weapons-systems. Identify the technical, operational, ethical, and legal considerations; identify what the engineering profession's collective response should be.
\end{exercise}

\begin{exercise}
Discuss the dual-use-technology export-control regime as an engineering responsibility. Identify the engineering decisions that determine whether a technology is appropriately controllable; identify the institutional context in which the engineering decisions are made.
\end{exercise}

